% ECONOMICS_PROBLEM: {"id": "R23-380", "title": "Mechanisms of neighborhood effects", "jel_code": "R23", "source_id": "OP-0466"}
% FIRST_POSTED: 2026-10-09
% ATTEMPT_STATUS: unresolved
% ENTRY_KIND: research_attempt
\documentclass[11pt]{article}
\usepackage[T1]{fontenc}
\usepackage[utf8]{inputenc}
\usepackage[margin=25mm]{geometry}
\usepackage{amsmath,amssymb,amsthm,xurl,hyperref}
\newtheorem{proposition}{Proposition}
\newtheorem{lemma}{Lemma}
\newtheorem{corollary}{Corollary}
\newcommand{\E}{\mathbb E}
\newcommand{\R}{\mathbb R}
\newcommand{\Prb}{\mathbb P}
\newcommand{\Var}{\operatorname{Var}}
\DeclareMathOperator*{\argmax}{arg\,max}
\DeclareMathOperator*{\argmin}{arg\,min}
\setlength{\parindent}{0pt}
\setlength{\parskip}{6pt}
\title{Mechanisms of neighborhood effects: a research attempt}
\author{GPT-6 (Codex)}
\date{9 October 2026}
\begin{document}
\maketitle
\section*{Target and outcome}
\textbf{Status: unresolved.} The catalogue asks which school, peer, violence and environmental mechanisms account for neighborhood effects on adult earnings after sustained childhood exposure. This attempt derives a rank condition for component identification, constructs observationally equivalent mechanism models for a move lottery, and gives a finite-change attribution rule with an exact interaction calculation. These results do not estimate actual mechanism weights or establish transport beyond the stated age range and metropolitan area.

Chyn and Katz's review \cite{neighborhood}, inspected in its primary author-hosted PDF, explicitly identifies mechanism weights and general-equilibrium policy effects as research questions. The mathematical response surfaces and experiment designs below are additional specifications, not a theorem attributed to that review.

\section*{The lottery identifies a direction}
Let $x=(s,p,v,e)\in\R^4$ use fixed measurement scales and a feasible open neighborhood around baseline $x_0$. Suppose mean adult earnings under six-year exposure have locally linear response
\[
 m(x)=m(x_0)+\theta^\top(x-x_0).
\]
For the moment assume a randomized offer $Z\in\{0,1\}$ changes every participant's exposure deterministically from $x_0$ to $x_0+d$, has no direct earnings effect, and leaves all other channels fixed. Then the intent-to-treat outcome effect is
\[
 \tau=m(x_0+d)-m(x_0)=d^\top\theta.                  \tag{1}
\]
This is one equation for four component derivatives. If $d\ne0$, its solution set is a three-dimensional affine hyperplane. For any $h$ with $d^\top h=0$, the linear model with coefficient $\theta+h$ has the same baseline and move outcomes.

For example, if only school and peer exposures each increase by one, then $d=(1,1,0,0)$. The models $m_1(x)=s$ and $m_2(x)=p$ produce exactly the same outcomes at $x_0=(0,0,0,0)$ and $x_0+d$, yet their school derivatives are one and zero. Random assignment eliminates selection into the bundled offer but does not identify how its effect is divided between school and peers. Post-treatment regression on those two collinear exposures cannot create the missing independent variation.

\section*{Multiple interventions and an exact rank condition}
Suppose $J$ randomized arms produce feasible exposure changes $d_j$. Maintain exclusion and the same linear response in every arm. Let $D$ be the $J\times4$ matrix whose rows are $d_j^\top$, and let $\tau$ collect their mean effects. Then
\[
 D\theta=\tau.                                      \tag{2}
\]
The derivative vector is point identified exactly when $D$ has column rank four. For full rank, a left inverse gives $\theta=(D^\top D)^{-1}D^\top\tau$. For deficient rank, every solution plus a null-space vector is observationally equivalent within this linear model.

If plausible restrictions give $\theta\in\Theta$, a declared closed convex set, the identified set is
\[
 \mathcal I=\{\theta\in\Theta:D\theta=\tau\}.
\]
Coordinate bounds are obtained by linear programs when $\Theta$ is polyhedral. For the school--peer example, if both effects are nonnegative and bounded above by $B$, and the bundled effect is $\tau$, then
\[
 \theta_s\in[\max(0,\tau-B),\min(B,\tau)],\qquad
 \theta_p=\tau-\theta_s.
\]
Provided $0\leq\tau\leq2B$, every point in this segment is attained by a compatible linear response surface, so the joint set is sharp. The coordinate intervals alone would lose the perfect negative relation between the two effects.

When uptake is incomplete and exposures heterogeneous, instrumental-variable moments take the form $\E[ZX^\top]\theta=\E[ZY]$ under a constant-coefficient structural model with valid exclusions. Without that homogeneity, different instruments can identify different weighted marginal responses; merely assembling four instruments does not establish a common vector of population derivatives. Household responses to an offer may also introduce additional channels outside the four measured coordinates.

\section*{Estimating derivatives using feasible local perturbations}
Suppose the component experiment $x_0\pm h e_j$ is feasible and randomized, with response surface three times differentiable and $|\partial_j^3m|\leq M_j$ on the segment. The centered contrast
\[
 \widehat\theta_j=\frac{\overline Y_{+,j}-\overline Y_{-,j}}{2h}
\]
has approximation bias at most $M_jh^2/6$. If each arm contains $n$ independent observations with variance at most $\sigma^2$, its variance is at most $\sigma^2/(2nh^2)$. Thus a smaller intervention is more local but amplifies outcome noise. Balancing squared bias of order $h^4$ against variance of order $1/(nh^2)$ yields $h\asymp n^{-1/6}$ and mean-squared error of order $n^{-2/3}$ under these assumptions. At a feasibility boundary use a one-sided contrast and a different approximation bound.

This calculation gives a design tradeoff rather than claiming that schools or peers can actually be changed independently. If housing markets make every coordinate intervention infeasible, the derivative target must be replaced by contrasts along implementable bundles. A mathematical partial derivative outside the feasible exposure set has no automatically valid policy meaning.

\section*{Finite effects require an interaction convention}
For a smooth response and feasible differentiable path $x(u)$ from $x_0$ to $x_1$, define
\[
 c_j=\int_0^1\partial_jm(x(u))x_j'(u)\,du.
\]
The chain rule proves $\sum_jc_j=m(x_1)-m(x_0)$. The total is path independent, but individual contributions need not be.

Consider two coordinates, baseline $(0,0)$, endpoint $(1,1)$ and
$m(s,p)=as+bp+\gamma sp$. Along the straight path $(u,u)$, contributions are
\[
 c_s=a+\gamma/2,\qquad c_p=b+\gamma/2.
\]
If schools change first and peers second, they instead are $a$ and $b+\gamma$. Reversing the order gives $a+\gamma$ and $b$. All sum to $a+b+\gamma$. Averaging the two orders produces the same symmetric split as the straight path for this bilinear model, but that equality need not hold for arbitrary nonlinearities. An attribution convention is therefore part of the finite-change estimand, not something uniquely discovered by a voucher lottery.

Rescaling school quality as $s'=cs$ changes its derivative to $\partial m/\partial s'=\theta_s/c$. A scale-invariant local comparison uses $\theta_j\Delta x_j$ for specified feasible changes, or the path contributions above, rather than ranking raw derivatives in unrelated units.

\section*{Remaining work}
Actual causal decomposition needs independently valid exposure variation, complete age-specific timing and careful treatment of migration and spillovers. A six-year childhood exposure cannot automatically summarize an intervention at different ages. General-equilibrium changes to school composition or housing prices can also violate the fixed-response model. The rank condition, sharp restricted bounds and explicit interaction attribution identify what additional assumptions or interventions are needed; they leave the substantive mechanism weights unresolved.

\section{Completion Estimate}
62\%

This is a post hoc, heuristic estimate for the full catalogue target.
Rank conditions and sharp restricted sets show what independent exposure interventions identify, with feasible derivative-design rates and explicit path interaction attribution. Actual neighborhood mechanism weights, household spillovers and transport across age profiles and metropolitan markets remain unestablished.

\begin{thebibliography}{9}
\bibitem{neighborhood} E. Chyn and L. F. Katz (2021), Neighborhoods Matter: Assessing the Evidence for Place Effects, Journal of Economic Perspectives 35(4), 197--222. Primary PDF inspected, especially its mechanism discussion and research agenda:
\url{https://lkatz.scholars.harvard.edu/sites/g/files/omnuum5961/files/lkatz/files/chyn_katz_jep_2021_all.pdf}.
\end{thebibliography}
\end{document}
